% !PRINTED 346
\weberpage{346}
If one applies the formula (1) to a real equation, then $\alpha^m + \beta^m + \gamma^m + \cdots$ is a real number; and if the absolutely greatest root $\alpha$ is real, then the $m$th root is also to be taken real. But if $m$ is an even number, then, in order to decide about the sign, one must further know whether $\alpha$ is positive or negative. If necessary, this is decided by substituting the approximate value found into the given equation.

Likewise, if one applies the formula (1) to an imaginary equation, one will have to decide which of the various $m$th roots gives the correct approximation.

The expressions
\[
s_m = \alpha^m + \beta^m + \gamma^m + \cdots
\]
can be computed especially simply when one takes for $m$ the successive powers of 2, that is $m = 2, 4, 8, 16 \ldots$, and they then yield very good results.

Now $s_2$ is the negative coefficient of the $(n-1)$th power of the unknown in the equation whose roots are the squares of the roots of $f(x)$. This equation, however, is easily formed in the following way.

In $f(x)$ collect the terms with even and with odd powers of $x$, and accordingly set
\[
f(x) = \varphi(x^2) + x\,\psi(x^2).
\]
Then
\[
f_1(x) = \varphi(x)^2 - x\,\psi(x)^2 = 0
\]
is the equation whose roots are the squares of the roots of $f(x)$. If one treats $f_1(x)$ in the same way, one obtains an equation whose roots are the fourth powers of the roots of $f(x)$, and so on.

The carrying out of this computation is very simple, and for the most part leads after a few steps to a good approximation.

We consider some examples.

First we take the equation
\weq{4}{x^3 - 2x - 2 = 0,}
which has one real positive and two imaginary roots. The real root is greater than 1,7, and since the product of all three roots is equal to 2 and $(1,7)^3 > 2$, the absolute value of the two imaginary roots is smaller than the real root. Hence the real root is the absolutely greatest, and the \gesperrt{Gräffe}-
%CONT>>
